\documentclass[10pt,a4paper]{article}
\usepackage{amsmath,amssymb,amsthm,graphicx}
\usepackage{amsmath,amssymb,amsthm,hyperref,xfrac}
\providecommand{\eg}{e.g.\ }
\providecommand{\ie}{i.e.\ }
\usepackage[scr=rsfs,cal=euler]{mathalfa}
\multlinegap0pt
\hyphenation{accu-rate applied assump-tion assump-tions decom-pos-ing defined esti-mate esti-mates exten-ded hypothe-ses para-me-ter para-me-ters remain-der required}
\newenvironment{enumeratei}
{\bgroup\def\theenumi{\roman{enumi}}\def\theenumii{\arabic{enumii}}\begin{enumerate}}
{\end{enumerate}\egroup}
\newenvironment{enumeratea}
{\bgroup\def\theenumi{\alph{enumi}}\begin{enumerate}}
{\end{enumerate}\egroup}

\newcommand{\Psfrac}[2]{(\sfrac{#1}{#2})}
\newcommand{\psfrac}[2]{\sfrac{(#1)}{#2}}
\newcommand{\spfrac}[2]{\sfrac{#1}{(#2)}}
\newcommand{\pspfrac}[2]{\sfrac{(#1)}{(#2)}}
\newcommand{\shalf}{\sfrac{1}{2}}

\newcommand{\bnorm}[1]{\bigl\|#1\bigr\|}
\newcommand{\Bnorm}[1]{\Bigl\|#1\Bigr\|}
\newcommand{\babs}[1]{\bigl|#1\bigr|}
\newcommand{\Babs}[1]{\Bigl|#1\Bigr|}

\AtBeginDocument{%
\DeclareMathOperator{\sinn}{sin}\let\sin\sinn
\DeclareMathOperator{\coos}{cos}\let\cos\coos
}

%\usepackage{enumitem}
\usepackage{physics}
%\usepackage[foot]{amsaddr}
\usepackage{bm}
\usepackage{tikz}
\usepackage{mathtools}
\usepackage{scalerel,stackengine}
\stackMath
\usetikzlibrary{arrows.meta}
\usepackage{subcaption}
%\usepackage[style=numeric]{biblatex}
%\addbibresource{main.bib}
%\setlength{\bibitemsep}{0.3\baselineskip}

\newenvironment{mycases}
{\begin{dcases*}}
{\end{dcases*}}
\theoremstyle{plain}
\newtheorem{theorem}{Theorem}[section]
\newtheorem{lemma}[theorem]{Lemma}
\newtheorem{definition}[theorem]{Definition}
\newtheorem{proposition}[theorem]{Proposition}
\newtheorem{corollary}[theorem]{Corollary}
\newtheorem*{corollary*}{Corollary}
\theoremstyle{remark}
\newtheorem{remark}[theorem]{Remark}
\counterwithout{figure}{section}
\counterwithout{equation}{section}

\newcommand{\br}[1]{\left(#1\right)}
\newcommand{\bbr}[1]{\bigl(#1 \bigr)}
\newcommand{\Bbr}[1]{\Bigl(#1 \Bigr)}
\newcommand{\bbbr}[1]{\biggl(#1 \biggr)}
\newcommand{\tpd}{\frac{1}{(2\pi)^{d}}}
\newcommand\pig[1]{\scalerel*[2.35pt]{\big#1}{%
\ensurestackMath{\addstackgap[1.2pt]{\big#1}}}}
\newcommand\pigl[1]{\mathopen{\pig{#1}}}
\newcommand\pigr[1]{\mathclose{\pig{#1}}}
\newcommand{\ak}{\abs{k}}
\newcommand{\ap}{\abs{p}}
\newcommand{\aal}{\abs{\alpha}}
\newcommand{\sbbr}[1]{\bigl[#1\bigr]}
\newcommand{\sBbr}[1]{\Bigl[#1\Bigr]}
\newcommand{\sbbbr}[1]{\biggl[#1\biggr]}
\newcommand{\sbr}[1]{\left[#1\right]}
\newcommand{\half}{\frac{1}{2}}
\newcommand{\p}{\partial}
\newcommand{\jb}[1]{\langle#1\rangle}
\newcommand{\bset}[1]{\bigl\{#1\bigr\}}
\newcommand{\Bset}[1]{\Bigl\{#1\Bigr\}}
\newcommand{\set}[1]{\left\{#1\right\}}
\newcommand{\mytextfrac}[2]{{\textstyle\frac{#1}{#2}}}
\newcommand{\qaq}{\quad\text{and}\quad}
\newcommand{\para}{\parallel}
\newcommand{\ep}{\varepsilon}
\newcommand{\R}{\mathbb{R}}
\newcommand{\C}{\mathbb{C}}
\newcommand{\al}{\alpha}
\newcommand{\intRd}{\int_{\R^d}}
\newcommand{\intRdk}{\int_{\R^d_k}}
\newcommand{\intRdp}{\int_{\R^d_p}}
\newcommand{\lam}{\lambda}
\newcommand{\ceils}[1]{\lceil#1\rceil}
\newcommand{\Qin}{Q_{\mathrm{in}}}
\newcommand{\hin}{h_{\mathrm{in}}}
\newcommand\reallywidehat[1]{%
\savestack{\tmpbox}{\stretchto{%
\scaleto{%
\scalerel*[\widthof{\ensuremath{#1}}]{\kern.pt\mathchar"0362\kern.1pt}%
{\rule{0ex}{\textheight}}%WIDTH-LIMITED CIRCUMFLEX
}{\textheight}%
}{2.4ex}}%
\stackon[-6.9pt]{#1}{\tmpbox}%
}
\title{Uniform decay of every Fourier mode of the density under free Schrödinger evolution}
\author{Marnie Smith}
\date{doi:10.5802/jep.353; CC BY4.0}
\begin{document}
\maketitle
\noindent\textit{Editorial note.} The counted statement is the author's labelled proposition `propfreephasemixing` (source lines 550-555); the proof environment begins on the next line after \textbackslash{}end\{proposition\} apart from one blank line, so the pairing is by adjacency.\par\medskip
\noindent\textit{Editorial note.} Five context spans are retained, carrying the author's toolbox of quantum operators, the commutation remark, the two Wigner-transform identities and the free Schrödinger equation; every label used inside the retained text is therefore defined inside the package.\par\medskip
\noindent\textit{Editorial note.} The retained preamble is the publisher layout adapter plus the native macro, operator and theorem block of the same pinned source.\par\medskip
\noindent\textit{Editorial note.} No retained line refers to a label, citation or figure outside the retained spans, so the package declares no external reference.\par\medskip
\noindent\textit{Editorial scope.} One original result of the pinned publisher article, counted once: uniform decay of the Fourier modes of the density under free Schrödinger evolution (author label propfreephasemixing).\par\medskip
\noindent\textit{Allowed standard prerequisites (editorial).}\begin{itemize}
\item Free Schrödinger evolution. Q\_FS(t) = e\textasciicircum{}\{i hbar t Delta / 2\} Q\_in e\textasciicircum{}\{-i hbar t Delta / 2\} is the solution of the free Schrödinger equation.. It defines the density whose Fourier modes are estimated, and its equation is retained as context.
\item Commutation of the quantum operators. The quantum derivatives and position operators commute in the four combinations recorded in the author's remark.. It allows the identities for the individual operators to be extended to the joint weight at source line 557.
\item Weighted Sobolev norm H\textasciicircum{}\{N0\}\_\{M,M\}. The norm controls the kernel up to order N0 in both variables with the weight M = ceil((d+1)/2).. It is the norm on the right-hand side of the estimate.
\end{itemize}\medskip
\begin{definition}[Quantum derivatives and weights]\label{toolboxdefinition}
The quantum analogues of spatial and velocity derivatives are
\begin{align*}
\bm{\nabla}_xQ&:=[\nabla,Q],\quad\hspace{0.045cm}\, (\bm{\nabla}_xQ)(x,y)=(\nabla_x+\nabla_y)Q(x,y),\\
\bm{\nabla}_\xi Q&:=[\mytextfrac{x}{i\hbar},Q],\quad \hspace{0.015cm}(\bm{\nabla}_\xi Q)(x,y)=(\mytextfrac{x-y}{i\hbar})Q(x,y).
\end{align*}
The quantum analogues of spatial and velocity weights are
\begin{alignat*}{2}
\bm{x}Q
&:=\mytextfrac{1}{2}(xQ+Qx),\quad\quad\,\,
(\bm{x}Q)(x,y)&&=(\mytextfrac{x+y}{2})Q(x,y),\\
\bm{\xi}Q&:=-\mytextfrac{i\hbar}{2}(\nabla Q+Q\nabla),\quad
(\bm{\xi}Q)(x,y)
&&=-\mytextfrac{i\hbar}{2}(\nabla_x-\nabla_y)Q(x,y).
\end{alignat*}
\end{definition}

\begin{remark}[Commutation of the quantum operators]\label{commutationremark}
The operators of Definition~\ref{toolboxdefinition} satisfy the commutation relations
\begin{align*}
[\bm{\nabla}_x,\bm{\nabla}_\xi]=[\bm{x},\bm{\xi}]=[\bm{\nabla}_x,\bm{\xi}]=[\bm{x},\bm{\nabla}_\xi]=0,
\end{align*}
in exact analogy with the classical operators $\nabla_x$, $\nabla_\xi$, $x$, $\xi$ acting on functions $f(x,\xi)$. Indeed, on the level of integral kernels, $\bm{x}$ and $\bm{\nabla}_x$ involve only the sum $x+y$, while~$\bm{\xi}$ and~$\bm{\nabla}_\xi$ involve only the difference $x-y$: for instance,
\begin{align*}
(\nabla_x-\nabla_y)\sbr{\br{\mytextfrac{x+y}{2}}Q(x,y)}=\br{\mytextfrac{x+y}{2}}(\nabla_x-\nabla_y)Q(x,y),
\end{align*}
which gives $[\bm{x},\bm{\xi}]=0$, and similarly $(\nabla_x+\nabla_y)(x-y)=0$ gives $[\bm{\nabla}_x,\bm{\nabla}_\xi]=0$. The remaining relations hold since $\bm{\nabla}_x$ and $\bm{\xi}$ are both constant-coefficient differential operators, and $\bm{x}$ and $\bm{\nabla}_\xi$ are both multiplication operators. In particular, products such as $\bm{\xi}^{\alpha}\bm{x}^{\beta}$ are independent of the order of the factors. By contrast, $[\bm{\nabla}_x,\bm{x}]\neq0$ and $[\bm{\nabla}_\xi,\bm{\xi}]\neq0$, again as in the classical case.
\end{remark}

\begin{lemma}\label{FTofWignertransform}
Suppose $Q\in\mathcal H^0_{M,M}(\R^d;\C)$. Then the Wigner transform of $Q$ satisfies, for all $k,\eta\in\R^d$,
\begin{align*}
\babs{\widehat {W[Q]}(k,\eta)}\lesssim_d\norm{Q}_{\mathcal H^{0}_{M,M}}.
\end{align*}
\end{lemma}

\begin{lemma}\label{operatorsWigner}Suppose $Q\in\mathcal L^2(\R^d;\C)$. Then the Wigner transform of $Q$ satisfies the following:\begin{alignat*}{2}
\nabla_xW[Q](x,\xi)&=W[\bm{\nabla}_x Q](x,\xi),\qquad ik\widehat{W[Q]}(k,\eta)&&=\reallywidehat{W[\bm{\nabla}_x Q]}(k,\eta), \\
\nabla_\xi W[Q](x,\xi)&=W[\bm{\nabla}_\xi Q](x,\xi),\qquad i\eta \widehat{W[Q]}(k,\eta)&&=\reallywidehat{W[\bm{\nabla}_\xi Q]}(k,\eta), \\
xW[Q](x,\xi)&=W[\bm{x}Q](x,\xi),\qquad i\nabla_k\widehat{W[Q]}(k,\eta)&&=\reallywidehat{W[\bm{x}Q]}(k,\eta), \\
\xi W[Q](x,\xi)&=W[\bm{\xi}Q](x,\xi),\qquad i\nabla_\eta\widehat{W[Q]}(k,\eta)&&=\reallywidehat{W[\bm{\xi}Q]}(k,\eta).
\end{alignat*}
\end{lemma}

\begin{align}\label{FS}
\begin{mycases}
i\hbar\p_tQ_{\mathrm{FS}}=\sBbr{-\frac{\hbar^2}{2}\Delta,Q_{\mathrm{FS}}},\\
Q_{\mathrm{FS}}(0)=\Qin,
\end{mycases}
\end{align}

\begin{proposition}[Phase mixing for the free Schrödinger equation] \label{propfreephasemixing}
Fix $N_0\geq0$. Then for any $t\geq0$ and $k\in\R^d$, the density of the solution $Q_{\mathrm{FS}}$ to \eqref{FS} satisfies
\begin{align*}
\abs{\widehat{\rho_{{Q_{\mathrm{FS}}}}}(t,k)}\lesssim_d\frac{ \norm{Q_{\mathrm{in}}}_{\mathcal H^{N_0}_{M,M}}}{\jb{k,kt}^{N_0}}.
\end{align*}
\end{proposition}

\begin{proof}
First recall that $\widehat{\rho_{Q_{\mathrm{FS}}}}(t,k)=\widehat{W[\Qin]}(k,kt)$. Since $\bm{\nabla}_x$ and $\bm{\nabla}_\xi$ commute (Remark~\ref{commutationremark}), the relations of Lemma~\ref{operatorsWigner} extend from the individual operators to the joint weight $\jb{\bm{\nabla}_x,\bm{\nabla}_\xi}^{N_0}$, whence
\begin{align*}
\jb{k,kt}^{N_0}\widehat{\rho_{Q_{\mathrm{FS}}}}(t,k)=\jb{k,\eta}^{N_0}\widehat{W[\Qin]}(k,\eta)\Big|_{\eta=kt}=\reallywidehat{W\sbr{\jb{\bm{\nabla}_x,\bm{\nabla}_{\xi}}^{N_0}Q_{\mathrm{in}}}}(k,kt).
\end{align*}
The result then follows from Lemma~\ref{FTofWignertransform}.
\end{proof}

\begin{thebibliography}{99}
\bibitem{r1} J. Bedrossian - “Nonlinear echoes and Landau damping with insufficient regularity”, Tunis. J. Math. 3 (2021) no. 1, p. 121-205
\bibitem{r2} J. Bedrossian, N. Masmoudi \& C. Mouhot - “Landau damping: paraproducts and Gevrey regularity”, Ann. PDE 2 (2016) no. 1, article ID 4, 71 pages
\bibitem{r3} J. Bedrossian, N. Masmoudi \& C. Mouhot - “Landau damping in finite regularity for unconfined systems with screened interactions”, Comm. Pure Appl. Math. 71 (2018) no. 3, p. 537-576
\bibitem{r4} J. Bedrossian, N. Masmoudi \& C. Mouhot - “Linearized wave-damping structure of Vlasov–Poisson in $\mathbb{R}^3$”, SIAM J. Math. Anal. 54 (2022) no. 4, p. 4379-4406
\bibitem{r5} A. Borie, S. Hadama \& J. Sabin - “Scattering for the positive density Hartree equation”, 2025
\bibitem{r6} A. Bove, G. Da Prato \& G. Fano - “An existence proof for the Hartree-Fock time-dependent problem with bounded two-body interaction”, Comm. Math. Phys. 37 (1974), p. 183-191
\bibitem{r7} A. Bove, G. Da Prato \& G. Fano - “On the Hartree-Fock time-dependent problem”, Comm. Math. Phys. 49 (1976) no. 1, p. 25-33
\bibitem{r8} T. Cazenave - Semilinear Schrödinger equations, Courant Lecture Notes in Math., vol. 10, American Mathematical Society, Providence, RI, 2003
\bibitem{r9} J. M. Chadam - “The time-dependent Hartree-Fock equations with Coulomb two-body interaction”, Comm. Math. Phys. 46 (1976) no. 2, p. 99-104
\bibitem{r10} T. Chen, Y. Hong \& N. Pavlović - “Global well-posedness of the NLS system for infinitely many fermions”, Arch. Rational Mech. Anal. 224 (2017) no. 1, p. 91-123
\bibitem{r11} T. Chen, Y. Hong \& N. Pavlović - “On the scattering problem for infinitely many fermions in dimensions $d\ge 3$ at positive temperature”, Ann. Inst. H. Poincaré C Anal. Non Linéaire 35 (2018) no. 2, p. 393-416
\bibitem{r12} J. J. Chong, L. Lafleche \& C. Saffirio - “On the $L^2$ rate of convergence in the limit from the Hartree to the Vlasov-Poisson equation”, J. Éc. polytech. Math. 10 (2023), p. 703-726
\bibitem{r13} C. Collot \& A.-S. de Suzzoni - “Stability of equilibria for a Hartree equation for random fields”, J. Math. Pures Appl. (9) 137 (2020), p. 70-100
\bibitem{r14} C. Collot \& A.-S. de Suzzoni - “Stability of steady states for Hartree and Schrödinger equations for infinitely many particles”, Ann. H. Lebesgue 5 (2022), p. 429-490
\bibitem{r15} E. Grenier, T. T. Nguyen \& I. Rodnianski - “Landau damping for analytic and Gevrey data”, Math. Res. Lett. 28 (2021) no. 6, p. 1679-1702
\bibitem{r16} S. Hadama - “Asymptotic stability of a wide class of stationary solutions for the Hartree and Schrödinger equations for infinitely many particles”, Ann. H. Lebesgue 8 (2025), p. 181-218
\bibitem{r17} A. D. Ionescu, B. Pausader, X. Wang \& K. Widmayer - “Nonlinear Landau damping for the Vlasov-Poisson system in $\mathbb{R}^3$: the Poisson equilibrium”, Ann. PDE 10 (2024) no. 1, article ID 2, 78 pages
\bibitem{r18} L. Landau - “On the vibrations of the electronic plasma”, Acad. Sci. USSR. J. Phys. 10 (1946), p. 25-34
\bibitem{r19} M. Lewin \& J. Sabin - “The Hartree equation for infinitely many particles, II: Dispersion and scattering in 2D”, Anal. PDE 7 (2014) no. 6, p. 1339-1363
\bibitem{r20} M. Lewin \& J. Sabin - “The Hartree equation for infinitely many particles I. Well-posedness theory”, Comm. Math. Phys. 334 (2015) no. 1, p. 117-170
\bibitem{r21} M. Lewin \& J. Sabin - “The Hartree and Vlasov equations at positive density”, Comm. Partial Differential Equations 45 (2020) no. 12, p. 1702-1754
\bibitem{r22} C. Mouhot \& C. Villani - “On Landau damping”, Acta Math. 207 (2011) no. 1, p. 29-201
\bibitem{r23} T. T. Nguyen \& C. You - “Modified scattering for long-range Hartree equations of infinite rank near vacuum”, SIAM J. Math. Anal. 57 (2025) no. 6, p. 6486-6497
\bibitem{r24} T. T. Nguyen \& C. You - “Plasmons for the Hartree equations with Coulomb interaction”, Probab. Math. Phys. 6 (2025) no. 3, p. 913-960
\bibitem{r25} E. M. Stein \& G. Weiss - “Interpolation of operators with change of measures”, Trans. Amer. Math. Soc. 87 (1958), p. 159-172
\bibitem{r26} C. You - “Phase mixing estimates for the nonlinear Hartree equation of infinite rank”, 2024
\end{thebibliography}
\end{document}
